\documentclass[10pt,a4paper]{amsart}
\usepackage[margin=19mm]{geometry}
\usepackage{amsmath,amssymb,mathtools}
\usepackage[scr=rsfs,cal=euler]{mathalfa}
\usepackage{xfrac}
\usepackage[unicode,hidelinks,bookmarks=false]{hyperref}
\newcommand{\ie}{i.e.\ }
\newcommand{\cf}{cf.\ }
\newcommand{\resp}{respectively\ }
\newcommand{\Subsection}{\subsection}
\providecommand{\nobreakdash}{\nobreak\hbox{-}}
\setlength{\emergencystretch}{2em}
\multlinegap0pt
\usepackage{amssymb}
\newtheorem{theorem}{Theorem}
\title{List chromatic numbers and singular compactness}
\author{Shimon Garti}
\date{Source: arXiv:2111.11826v4 (CC BY 4.0)}
\begin{document}
\maketitle

\noindent\emph{Source and licence.} List chromatic numbers and singular compactness. Version arXiv:2111.11826v4 (CC BY 4.0), distributed by arXiv under the Creative Commons Attribution 4.0 International licence, http://creativecommons.org/licenses/by/4.0/, as recorded in the arXiv licence metadata for this exact version and shown on the abstract page https://arxiv.org/abs/2111.11826v4. Full licence notice: \texttt{licenses/arxiv-2111.11826-CC-BY-4.0.txt}.\par\smallskip

\noindent\textit{Editorial scope.} Counted unit: the article's theorem thmgrplist (source0.tex lines 349--351). The article's complete original proof of it is delivered with it (source0.tex lines 353--393, one \texttt{proof} environment of 41 lines). No mathematics is written, repaired or re-derived: the statement and the proof are the article's own lines, in the article's own notation, and no retained line is substituted or re-flowed.  No further source-local context is needed: every cross-reference of every retained line resolves inside the two delivered spans.  Nothing was omitted from any retained span, so the package carries no omission record.  No retained line contains a citation, so the wrapper carries no bibliography and no bibliography entry.  No retained line contains a figure environment, an image-inclusion command or drawing code, so no image is delivered and the package declares no asset dependency.  Editorial formatting only, in addition: an AMS \texttt{amsart} preamble that restates the article's own declarations and macros used by the retained lines, namely the environments \texttt{theorem} on separate counters, together with the standard packages those lines need.  The retained environments print in this document as Theorem 1 in order of appearance, whereas the article prints the same items with the numbering of its own full document.  The complete native source of the article as distributed in the arXiv e-print for this exact version is bundled unchanged as \texttt{sources/118787/source0.tex}.
\begin{theorem}
\label{thmgrplist} ${\rm GRP}({\rm List})$ implies SCH.
\end{theorem}

\begin{proof}
Suppose that $\mu$ is a strong limit cardinal, either regular or singular.
Let $G$ be a bipartite graph with a bipartition of $V^G$ into $A$ and $B$ such that $|A|=\mu$ and $|B|=\lambda<2^\mu$.
Let us prove that ${\rm List}(G)\leq\mu$.

Firstly, let $A=\{v_\alpha:\alpha\in\mu\}$ and $B=\{w_\beta:\beta\in\lambda\}$.
Secondly, let $F$ be any $\mu$-assignment defined on $A\cup B$.
By induction on $\alpha\in\mu$ we choose $\eta_f\in F(v_\alpha)$ for every function $f\in{}^\alpha 2$ such that $f\neq g\Rightarrow\eta_f\neq\eta_g$ whenever $\beta\leq\alpha$ and $g\in{}^\beta 2$.
The choice is possible since $\mu$ is strong limit.
Specifically, if $\alpha\in\mu$ let $\theta=\bigcup\{2^\beta:\beta\leq\alpha\}$ so $\theta<\mu$.
Since $|F(v_\alpha)|=\mu$ one can choose a distinct element $\eta_f\in F(v_\alpha)$ for every $f\in\bigcup\{{}^\beta 2:\beta\leq\alpha\}$.

Suppose that $h\in{}^\mu 2$ and let $T_h=\{\eta_{h\upharpoonright\alpha}:\alpha\in\mu\}$.
By our choice if $\alpha\neq\beta$ then $\eta_{h\upharpoonright\alpha}\neq\eta_{h\upharpoonright\beta}$ and hence $|T_h|=\mu$.
On the other hand if $g\neq h$ then $|T_g\cap T_h|<\mu$ since $g\upharpoonright\beta\neq h\upharpoonright\beta$ from some $\beta$ onwards.
It follows that if $\beta\in\lambda$ then there is at most one function $h\in{}^\mu 2$ for which $F(w_\beta)\subseteq T_h$.
Since $\lambda<2^\mu$ one can choose $h\in{}^\mu 2$ such that $\forall\beta\in\lambda,\neg(F(w_\beta)\subseteq T_h)$.
Define $f(v_\alpha)=\eta_{h\upharpoonright\alpha}$ for every $\alpha\in\mu$ and choose $f(w_\beta)\in F(w_\beta)-T_h$ for every $\beta\in\lambda$.
By the above considerations, $f$ is a good coloring of $G$ and hence ${\rm List}(G)\leq\mu$.

Consider now the complete bipartite graph $K_{\mu,2^\mu}$.
Let $A\cup B$ be a bipartition of this graph with $A=\{a_\alpha:\alpha\in\mu\}$ and $B=\{b_\beta:\beta\in 2^\mu\}$.
Let us show that ${\rm List}(K_{\mu,2^\mu})>\mu$.

As a first step choose $\mathcal{A}=(A_\alpha:\alpha\in\mu)$ such that $|A_\alpha|=\mu$ for every $\alpha\in\mu$ and $\alpha_0\neq\alpha_1\Rightarrow A_{\alpha_0}\cap A_{\alpha_1}=\varnothing$.
Now let $E=\prod\{A_\alpha:\alpha\in\mu\}$, so $|E|=\mu^\mu=2^\mu$.
Define a $\mu$-assignment $F:A\cup B\rightarrow[{\rm Ord}]^\mu$ as follows.
For every $\alpha\in\mu$ let $F(a_\alpha)=A_\alpha$.
Fix a bijection $h:B\rightarrow E$, say $h(b_\beta)=g_\beta\in E$ for every $\beta\in 2^\mu$, and define $F(b_\beta)=\{g_\beta(\alpha):\alpha\in\mu\}$.
Notice that $|F(b_\beta)|=\mu$ by the disjointness of the $A_\alpha$s.

Suppose that $f$ is a coloring defined on $A\cup B$ so that $f(v)\in F(v)$ for every $v\in A\cup B$.
Since $f\upharpoonright{A}\in E$ there exists a unique $\beta\in 2^\mu$ such that $f\upharpoonright{A}=g_\beta$.
It follows that $f(b_\beta)\in F(b_\beta)=\{g_\beta(\alpha):\alpha\in\mu\}=\{f(a_\alpha):\alpha\in\mu\}$, so for some $a_\alpha\in A$ we have $f(b_\beta)=f(a_\alpha)$.
This fact shows that $f$ is not a good coloring since every element of $A$ is connected with every element of $B$, hence ${\rm List}(K_{\mu,2^\mu})>\mu$ as desired.

Assume now that ${\rm GRP}({\rm List})$ holds and $\mu$ is a strong limit cardinal.
If $2^\mu>\mu^+$ then $G=K_{\mu,2^\mu}$ is a graph of size greater than $\mu^+$ and ${\rm List}(G)>\mu$.
But if $H\leq G$ and $|H|=\mu^+$ then $H$ is a subgraph of $K_{\mu,\mu^+}$ and hence ${\rm List}(H)\leq\mu$ contradicting ${\rm GRP}({\rm List})$.
We conclude, therefore, that $2^\mu=\mu^+$ as required.
\end{proof}

\end{document}
